%% Chapter 12 ----------------------------------------------------------------
\chapter{Multi-Station Comparison}
\label{ch:comparison}
\index{Multi-Station Comparison}

The Multi-Station Comparison workspace displays dynamic spectra from several stations in
vertically stacked panels with a shared time axis, so that the same event can be compared
across stations, receivers and frequency ranges (Figure~\ref{fig:comparison}). It is opened
from \menu{View > Multi-Station Comparison...} or with the \gui{Compare} button of the
downloader (Chapter~\ref{ch:downloader}).

\screenshot{downloader/multi_station_comparison}{The Multi-Station Comparison workspace with four stations observing the same event.}{fig:comparison}{Multi-Station Comparison window.}

The workspace follows the rendering mode of the application: in Modern mode the panels are
hardware-accelerated when available, and in Classic mode they are drawn with Matplotlib\index{Matplotlib}.

\section{Stations and files}

The \gui{Stations / Files} group lists the loaded files as \emph{station -- date -- file
name}. \gui{Add Files...} adds FITS files from disk, \gui{Remove} removes the selected entry
and \gui{Clear} removes all entries. Files that can be combined are merged automatically
before display: each station's files are combined in time, in frequency, or as a complete
time~\(\times\)~focus-code grid (Section~\ref{sec:combine-rules}), so a mixed-station
selection yields one panel per station. Hovering over an entry shows any warning raised
while reading it.

\section{Display controls}

\begin{table}[!htbp]
  \centering
  \caption{Display controls of the Multi-Station Comparison.}
  \label{tab:comparison-controls}
  \small
  \begin{tabularx}{\linewidth}{@{}P{0.15\linewidth}P{0.2\linewidth}L@{}}
    \toprule
    \thead{Group} & \thead{Control} & \thead{Function} \\
    \midrule
    View & Time alignment\index{Multi-Station Comparison!time alignment} & \gui{UT clock} aligns the panels on Universal Time; \gui{Seconds from file start} aligns each panel on its own start. UT alignment requires \texttt{TIME-OBS}\index{TIME-OBS keyword} in every file; otherwise the workspace switches to seconds and says so. \\
    Appearance & Units & \gui{Digits} or \gui{dB}. \\
               & Colormap & Color map of all panels. \\
               & Color scale\index{Multi-Station Comparison!color scale} & \gui{Shared scale} (one scale for all panels), \gui{Per-station auto scale}, or \gui{Manual scale} with \gui{Low} and \gui{High} limits. \\
    Noise Reduction & Target & \gui{All panels} or one panel. \\
                    & Method & \gui{None}, \gui{Mean background}, \gui{Median background}, \gui{Robust background} or \gui{Noise clipping}. \\
                    & Colormap & Color map of the noise-reduced panels. \\
                    & Thresholds & With \gui{Noise clipping}, \gui{Lower threshold} and \gui{Upper threshold} sliders and a \gui{Logarithmic Threshold Scale}\index{display thresholds!logarithmic scale} option, as in the main window. \\
    \bottomrule
  \end{tabularx}
\end{table}

\section{Display range and view configurations}

The \gui{Actions} group contains:
\begin{description}
  \item[Set Range...] enter a shared time and frequency range for all panels.
  \item[Reset Range] return to the full range of the data.
  \item[Load View Config...] apply a \file{.efaview.json}\index{efaview.json@.efaview.json file} view configuration
    (Section~\ref{sec:view-config}). The visual settings are applied; when the files differ
    in date, set UT bounds with \gui{Set Range...}.
  \item[Export...] export the comparison (Section~\ref{sec:comparison-export}).
\end{description}

\section{Ruler measurement}

In the \gui{Measurement} group, \gui{Start Ruler}\index{Multi-Station Comparison!ruler} activates a ruler: click two points in one
panel to measure the duration, frequency drift and slope, shown under \gui{Readout}.
\gui{Clear} removes the measurement.

\section{Exporting a comparison}
\label{sec:comparison-export}

\gui{Export...} opens the \gui{Export Comparison} dialog (Figure~\ref{fig:comparison-export}):
\begin{description}
  \item[Visible View] exports the comparison exactly as it appears in the window.
  \item[Grid / Table Layout] exports a compact publication-style matrix of all panels, with
    a shared time axis, each station's native frequency range and a white background. Its
    options are the number of \gui{Columns} (\gui{Auto} or a fixed number), an
    \gui{Overall title}, and the \gui{Raster resolution} (DPI) for bitmap formats.
\end{description}
The figure can be saved as PNG, PDF, EPS, SVG or TIFF.

\screenshot[0.5\linewidth]{dialogs/comparison_export_dialog}{The \gui{Export Comparison} dialog.}{fig:comparison-export}{The Export Comparison dialog with Grid / Table Layout selected, showing the Columns, Overall title and Raster resolution fields and the explanatory note.}
